\chapter{\texorpdfstring{Optimization}{5. Optimization}}

\section{\texorpdfstring{Monotone Predicate Search}{5.1 Monotone Predicate Search}}
\setcounter{section}{1}
\setcounter{subsection}{0}
\subsection{Binary Search}
Binary search can be generally used to find the input value corresponding to any output value of a monotonic (strictly non-increasing or strictly non-decreasing) function in $O(\log n)$ time with respect to the domain size. This is a special case of finding the exact point at which any given monotonic Boolean function changes from true to false or vice versa. Each step evaluates the midpoint of the current interval and keeps the half that still contains the transition point, halving the search space until it closes on the answer. Unlike searching through an array, discrete binary search is not restricted by available memory, making it useful for handling infinitely large search spaces such as real number intervals.

\begin{compactitem}
  \item \inlinecode{binary\_search\_first\_true(lo, hi, pred)} takes signed integer boundaries for the search space $[{\inlinecodemath{lo}}, {\inlinecodemath{hi}})$ (i.e. including \inlinecode{lo}, but excluding \inlinecode{hi}) and returns the smallest integer \inlinecode{k} in $[{\inlinecodemath{lo}}, {\inlinecodemath{hi}})$ for which the predicate \inlinecode{pred(k)} tests true. If \inlinecode{pred(k)} tests false for the entire input range, then \inlinecode{hi} is returned. The caller must ensure \inlinecode{pred} is monotonic on the input range, i.e. returning all false for some (possibly empty) prefix, followed by all true in some (possibly empty) suffix. E.g., patterns \inlinecode{01}, \inlinecode{00}, and \inlinecode{11} are allowed, but \inlinecode{10} is disallowed.
  \item \inlinecode{binary\_search\_last\_true(lo, hi, pred)} takes signed integer boundaries for the search space $[{\inlinecodemath{lo}}, {\inlinecodemath{hi}})$ (i.e. including \inlinecode{lo}, but excluding \inlinecode{hi}) and returns the largest integer \inlinecode{k} in $[{\inlinecodemath{lo}}, {\inlinecodemath{hi}})$ for which the predicate \inlinecode{pred(k)} tests true. If \inlinecode{pred(k)} tests false for the entire input range, then the original \inlinecode{lo - 1} is returned, so \inlinecode{lo} must exceed the minimum representable value. The caller must ensure \inlinecode{pred} is monotonic on the input range, i.e. returning all true for some (possibly empty) prefix, followed by all false in some (possibly empty) suffix. E.g., patterns \inlinecode{10}, \inlinecode{00}, and \inlinecode{11} are allowed, but \inlinecode{01} is disallowed.
  \item \inlinecode{fbinary\_search(lo, hi, pred)} is the equivalent of \inlinecode{binary\_search\_first\_true()} on floating point predicates. Since any interval of real numbers is dense, the exact target cannot be found due to floating point error. Instead, a value that is very close to the border between false and true is returned. The precision of the answer depends on the number of repetitions the function performs. Since each repetition bisects the search space, the absolute error of the answer is $1/(2^n)$ times the distance between \inlinecode{lo} and \inlinecode{hi} after $n$ repetitions. Although the error can be controlled by looping until the distance shrinks to an arbitrary epsilon, it is simpler to let the loop run for a desired number of iterations until floating point arithmetic breaks down. $100$ iterations is usually sufficient, since the search space will be reduced to $2^{-100}$ (roughly $10^{-30}$) times its original size. This implementation can be modified to find the ``last true'' point by simply interchanging the assignments of \inlinecode{lo} and \inlinecode{hi} in the if-else statements.
\end{compactitem}

\Needspace*{3\baselineskip}
\textbf{Time Complexity}:
\begin{compactitem}
  \item $O(\log n)$ calls will be made to \inlinecode{pred()} in \inlinecode{binary\_search\_first\_true()} and \inlinecode{binary\_search\_last\_true()}, where $n$ is the distance between \inlinecode{lo} and \inlinecode{hi}.
  \item $O(n)$ calls will be made to \inlinecode{pred()} in \inlinecode{fbinary\_search()}, where $n$ is the number of iterations performed.
\end{compactitem}

\textbf{Space Complexity}: $O(1)$ auxiliary for all operations.

\begin{lstlisting}[language=C++]
template<typename Int, typename Pred>
Int binary_search_first_true(Int lo, Int hi, Pred pred) {  // 000[1]11
  while (lo < hi) {
    Int mid = lo + (hi - lo) / 2;
    if (pred(mid)) {
      hi = mid;
    } else {
      lo = mid + 1;
    }
  }
  return lo;  // hi if all false
}

template<typename Int, typename Pred>
Int binary_search_last_true(Int lo, Int hi, Pred pred) {  // 11[1]000
  while (lo < hi) {
    Int mid = lo + (hi - lo) / 2;
    if (pred(mid)) {
      lo = mid + 1;
    } else {
      hi = mid;
    }
  }
  return lo - 1;
}

template<typename Pred>
double fbinary_search(double lo, double hi, Pred pred) {  // 000[1]11
  double mid;
  for (int i = 0; i < 100; i++) {
    mid = (lo + hi) / 2.0;
    if (pred(mid)) {
      hi = mid;
    } else {
      lo = mid;
    }
  }
  return lo;
}
\end{lstlisting}
\Needspace*{4\baselineskip}
\textbf{Example Usage}\par
\begin{lstlisting}[language=C++]
#include <cassert>
#include <cmath>

int main() {
  assert(binary_search_first_true(0, 7, [](int x) { return x >= 3; }) == 3);
  assert(binary_search_first_true(0, 7, [](int x) { return true; }) == 0);
  assert(binary_search_first_true(0, 7, [](int x) { return false; }) == 7);
  assert(binary_search_first_true(4, 4, [](int x) { return x >= 4; }) == 4);

  assert(binary_search_last_true(0, 7, [](int x) { return x <= 5; }) == 5);
  assert(binary_search_last_true(0, 7, [](int x) { return true; }) == 6);
  assert(binary_search_last_true(0, 7, [](int x) { return false; }) == -1);
  assert(binary_search_last_true(4, 4, [](int x) { return x <= 4; }) == 3);

  double res = fbinary_search(-10.0, 10.0, [](double x) { return x >= 1.2345; });
  assert(fabs(res - 1.2345) < 1e-15);
  return 0;
}
\end{lstlisting}
\subsection{Exponential Search}
Finds the first value where a monotone Boolean predicate becomes true when an upper bound is not known in advance. Exponential search first grows the search range by powers of two until it brackets the transition point, then finishes with ordinary binary search.

This is useful for answer-search problems over unbounded or very large integer domains. The predicate must eventually become true, or the search may overflow or run forever unless an explicit limit is added.

\begin{compactitem}
  \item \inlinecode{exponential\_search\_first\_true(lo, pred)} returns the smallest integer \inlinecode{x} greater than or equal to \inlinecode{lo} such that \inlinecode{pred(x)} is true.
\end{compactitem}

\textbf{Time Complexity}: $O(\log n)$ calls to \inlinecode{pred()}, where $n$ is the distance from \inlinecode{lo} to the first true value.

\textbf{Space Complexity}: $O(1)$ auxiliary.

\begin{lstlisting}[language=C++]
#include <cstdint>

template<typename Int, typename Pred>
Int exponential_search_first_true(Int lo, Pred pred) {  // 000[1]11
  if (pred(lo)) {
    return lo;
  }
  Int step = 1;
  while (!pred(lo + step)) {
    step *= 2;
  }
  Int hi = lo + step;
  lo += step / 2;
  while (lo < hi) {
    Int mid = lo + (hi - lo) / 2;
    if (pred(mid)) {
      hi = mid;
    } else {
      lo = mid + 1;
    }
  }
  return lo;
}
\end{lstlisting}
\Needspace*{4\baselineskip}
\textbf{Example Usage}\par
\begin{lstlisting}[language=C++]
#include <cassert>

int main() {
  auto at_least_1000 = [](int x) { return x >= 1000; };
  assert(exponential_search_first_true(1000, at_least_1000) == 1000);
  assert(exponential_search_first_true(0, at_least_1000) == 1000);
  assert(exponential_search_first_true(5, at_least_1000) == 1000);
  auto square_large_enough = [](int64_t x) { return x * x >= 123456789LL; };
  assert(exponential_search_first_true(0LL, square_large_enough) == 11112);
  return 0;
}
\end{lstlisting}
\subsection{Local Minimum (Binary Search)}
Finds a non-strict local minimum in an array, that is, an index $i$ whose value is no greater than the values at its existing neighbors. Unlike ternary search, this does not require the array to be unimodal, but the result is only a local minimum and need not be the global minimum.

Compare the middle element with its right neighbor. If the sequence slopes downward, the right half contains a local minimum; otherwise, the middle element or the left half contains one. Each comparison therefore discards half of the remaining indices.

\begin{compactitem}
  \item \inlinecode{local\_min\_index(a)} returns the index of a non-strict local minimum in the nonempty array \inlinecode{a}. Either endpoint may be returned when it is no greater than its sole neighbor.
\end{compactitem}

\textbf{Time Complexity}: $O(\log n)$ comparisons, where $n$ is the size of \inlinecode{a}.

\textbf{Space Complexity}: $O(1)$ auxiliary.

\begin{lstlisting}[language=C++]
#include <cassert>
#include <vector>

template<typename T>
int local_min_index(const std::vector<T> &a) {
  assert(!a.empty());
  int lo = 0, hi = static_cast<int>(a.size()) - 1;
  while (lo < hi) {
    int mid = lo + (hi - lo) / 2;
    if (a[mid] > a[mid + 1]) {
      lo = mid + 1;
    } else {
      hi = mid;
    }
  }
  return lo;
}
\end{lstlisting}
\Needspace*{4\baselineskip}
\textbf{Example Usage}\par
\begin{lstlisting}[language=C++]
int main() {
  assert(local_min_index(std::vector<int>{9, 7, 3, 4, 8}) == 2);
  assert(local_min_index(std::vector<int>{1, 2, 3, 4}) == 0);

  std::vector<int> duplicates{4, 2, 2, 3};
  int i = local_min_index(duplicates);
  assert((i == 1 || i == 2) && duplicates[i] == 2);
  return 0;
}
\end{lstlisting}

\section{\texorpdfstring{Unimodal and Continuous Search}{5.2 Unimodal and Continuous Search}}
\setcounter{section}{2}
\setcounter{subsection}{0}
\subsection{Ternary and Golden Section Search}
Given a unimodal function $f$, find the position of its global minimum or maximum: a value $x$ such that the function strictly decreases up to $x$ and strictly increases after it (for a minimum), or the reverse (for a maximum). All three searches below exploit the same idea: evaluate two interior points, and the comparison reveals one side of the interval that cannot contain the optimum, so it can be discarded. They differ only in how the probes are placed and whether the domain is continuous or discrete. Each maximum routine is the minimum routine applied to the negated function, since the location of a maximum of $f$ is the location of a minimum of $-f$.

Ternary search splits the interval into thirds. It is the simplest to reason about, but it recomputes both probe points each step, spending two function evaluations per iteration while the interval shrinks to two-thirds.

Golden-section search places the probes at the golden ratio so that one probe is reused as an interior point of the next, smaller interval. It therefore performs only one new evaluation per iteration while shrinking the interval by a factor of $1/\varphi$ (about 0.618), and reaches a given accuracy in roughly 2.4 times fewer evaluations than ternary search. Prefer it when the function is expensive to evaluate.

The discrete version operates on an integer domain, narrowing by thirds until at most three candidates remain and then scanning them. It returns the index of the optimum rather than a real coordinate. As with ternary search on reals, the function must be strictly unimodal.

\begin{compactitem}
  \item \inlinecode{ternary\_search\_min(lo, hi, f, eps = 1e-12)} and \inlinecode{ternary\_search\_max(lo, hi, f, eps = 1e-12)} return a \inlinecode{double} \inlinecode{x} in $[{\inlinecodemath{lo}}, {\inlinecodemath{hi}}]$ optimizing the continuous unimodal function \inlinecode{f}, to within the optional absolute error \inlinecode{eps}.
  \item \inlinecode{golden\_section\_min(lo, hi, f, eps = 1e-12)} and \inlinecode{golden\_section\_max(lo, hi, f, eps = 1e-12)} do the same with one function evaluation per iteration.
  \item \inlinecode{discrete\_ternary\_min(lo, hi, f)} and \inlinecode{discrete\_ternary\_max(lo, hi, f)} return the integer index in range $[{\inlinecodemath{lo}}, {\inlinecodemath{hi}}]$ optimizing the unimodal function \inlinecode{f}, breaking ties toward the smaller index.
\end{compactitem}

\Needspace*{3\baselineskip}
\textbf{Time Complexity}:
\begin{compactitem}
  \item $O(\log\left(n / {\inlinecodemath{eps}}\right))$ calls to \inlinecode{f()} for the continuous searches, where $n$ is the distance between \inlinecode{lo} and \inlinecode{hi}. Golden-section makes about half as many per iteration as ternary search.
  \item $O(\log n)$ calls to \inlinecode{f()} for the discrete searches.
\end{compactitem}

\textbf{Space Complexity}: $O(1)$ auxiliary for all operations.

\begin{lstlisting}[language=C++]
template<typename Fn>
double ternary_search_min(double lo, double hi, Fn f, double eps = 1e-12) {
  while (hi - lo > eps) {
    double m1 = lo + (hi - lo) / 3, m2 = hi - (hi - lo) / 3;
    if (f(m1) < f(m2)) {
      hi = m2;
    } else {
      lo = m1;
    }
  }
  return (lo + hi) / 2;
}

template<typename Fn>
double ternary_search_max(double lo, double hi, Fn f, double eps = 1e-12) {
  return ternary_search_min(lo, hi, [&](double x) { return -f(x); }, eps);
}

template<typename Fn>
double golden_section_min(double lo, double hi, Fn f, double eps = 1e-12) {
  const double r = 0.6180339887498949;  // 1 / phi = (sqrt(5) - 1) / 2.
  double m1 = hi - r * (hi - lo), m2 = lo + r * (hi - lo);
  double f1 = f(m1), f2 = f(m2);
  while (hi - lo > eps) {
    if (f1 < f2) {
      hi = m2;
      m2 = m1;
      f2 = f1;
      m1 = hi - r * (hi - lo);
      f1 = f(m1);
    } else {
      lo = m1;
      m1 = m2;
      f1 = f2;
      m2 = lo + r * (hi - lo);
      f2 = f(m2);
    }
  }
  return (lo + hi) / 2;
}

template<typename Fn>
double golden_section_max(double lo, double hi, Fn f, double eps = 1e-12) {
  return golden_section_min(lo, hi, [&](double x) { return -f(x); }, eps);
}

template<typename Int, typename Fn>
Int discrete_ternary_min(Int lo, Int hi, Fn f) {
  while (hi - lo > 2) {
    Int m1 = lo + (hi - lo) / 3, m2 = hi - (hi - lo) / 3;
    if (f(m1) < f(m2)) {
      hi = m2 - 1;
    } else {
      lo = m1 + 1;
    }
  }
  Int best = lo;
  auto best_value = f(best);
  for (Int k = lo + 1; k <= hi; k++) {
    auto value = f(k);
    if (value < best_value) {
      best = k;
      best_value = value;
    }
  }
  return best;
}

template<typename Int, typename Fn>
Int discrete_ternary_max(Int lo, Int hi, Fn f) {
  return discrete_ternary_min(lo, hi, [&](Int x) { return -f(x); });
}
\end{lstlisting}
\Needspace*{4\baselineskip}
\textbf{Example Usage}\par
\begin{lstlisting}[language=C++]
#include <cassert>
#include <cmath>
#include <vector>
using namespace std;

bool EQ(double a, double b) {
  return fabs(a - b) < 1e-6;
}

int main() {
  // Parabola opening up with vertex at x = -2.
  auto up = [](double x) { return 3 * x * x + 12 * x - 12; };
  // Parabola opening down with vertex at x = 2/19.
  auto down = [](double x) { return -5.7 * x * x + 1.2 * x + 88; };

  assert(EQ(ternary_search_min(-1000, 1000, up), -2));
  assert(EQ(golden_section_min(-1000, 1000, up), -2));
  assert(EQ(ternary_search_max(-1000, 1000, down), 2.0 / 19));
  assert(EQ(golden_section_max(-1000, 1000, down), 2.0 / 19));

  // Discrete unimodal arrays: a valley and a peak.
  vector<int> valley{5, 3, 1, 2, 4, 6};
  vector<int> peak{1, 3, 7, 6, 2};
  auto v = [&](int i) { return valley[i]; };
  auto p = [&](int i) { return peak[i]; };
  assert(discrete_ternary_min(0, static_cast<int>(valley.size()) - 1, v) == 2);  // Value 1.
  assert(discrete_ternary_max(0, static_cast<int>(peak.size()) - 1, p) == 2);    // Value 7.
  return 0;
}
\end{lstlisting}
\subsection{2D Hill Climbing}
Given a continuous geometric objective $f(x, y)$ and a (possibly arbitrary) starting guess $(x_0, y_0)$, search for a small value using two-dimensional hill climbing.

The heuristic starts at the guess and evaluates one step in each of eight evenly spaced directions. It moves to the best improving neighbor, checks all eight directions again, and reduces the step size when no neighbor improves the answer. This repeats until the step size falls below the chosen threshold. Success depends heavily on the behavior of $f$ and the initial guess, so the result is not guaranteed to be the global minimum.

\begin{compactitem}
  \item \inlinecode{hill\_climb\_min(f, x0, y0, \&best\_x, \&best\_y, step\_min = 1e-9, step\_max = 1e6)} returns an approximate minimum value of function \inlinecode{f} reached by hill climbing from the starting guess $({\inlinecodemath{x0}}, {\inlinecodemath{y0}})$. If either optional pointer \inlinecode{best\_x} or \inlinecode{best\_y} is supplied, the corresponding coordinate of the point attaining the returned value is stored through it. The search starts with step size \inlinecode{step\_max} and stops below \inlinecode{step\_min}.
\end{compactitem}

\textbf{Time Complexity}: $O(k + s)$ calls to \inlinecode{f()}, where $k$ is the number of times the step size is reduced and $s$ is the total number of improving moves accepted.

\textbf{Space Complexity}: $O(1)$ auxiliary.

\begin{lstlisting}[language=C++]
template<typename Fn>
double hill_climb_min(
    Fn f, double x0, double y0, double *best_x = nullptr, double *best_y = nullptr,
    const double step_min = 1e-9, const double step_max = 1e6
) {
  static const double inv_sqrt2 = 0.7071067811865476;
  static const double dx[] = {1, inv_sqrt2, 0, -inv_sqrt2, -1, -inv_sqrt2, 0, inv_sqrt2};
  static const double dy[] = {0, inv_sqrt2, 1, inv_sqrt2, 0, -inv_sqrt2, -1, -inv_sqrt2};
  double x = x0, y = y0, res = f(x0, y0);
  for (double step = step_max; step > step_min;) {
    double next_value = res, next_x = x, next_y = y;
    bool found = false;
    for (int i = 0; i < 8; i++) {
      double x2 = x + step * dx[i], y2 = y + step * dy[i];
      double value = f(x2, y2);
      if (value < next_value) {
        next_x = x2;
        next_y = y2;
        next_value = value;
        found = true;
      }
    }
    if (!found) {
      step /= 2.0;
    } else {
      x = next_x;
      y = next_y;
      res = next_value;
    }
  }
  if (best_x != nullptr) {
    *best_x = x;
  }
  if (best_y != nullptr) {
    *best_y = y;
  }
  return res;
}
\end{lstlisting}
\Needspace*{4\baselineskip}
\textbf{Example Usage}\par
\begin{lstlisting}[language=C++]
#include <cassert>
#include <cmath>

bool EQ(double a, double b) {
  return fabs(a - b) < 1e-9;
}

// Paraboloid with global minimum at f(2, 3) = 0.
double f(double x, double y) {
  return (x - 2) * (x - 2) + (y - 3) * (y - 3);
}

int main() {
  double x, y;
  assert(EQ(hill_climb_min(f, 0, 0, &x, &y), 0));
  assert(EQ(x, 2) && EQ(y, 3));
  return 0;
}
\end{lstlisting}
\subsection{Simulated Annealing}
Approximately minimizes an energy function $E(s)$ over an arbitrary state space using simulated annealing. This randomized heuristic is useful when the states may be discrete, the energy has many local minima, and gradients or stronger structural properties are unavailable.

Starting from an initial state $s$, each iteration generates a random neighboring state $s'$. An improving move is always accepted, while a worsening move is accepted with probability $\exp(-(E(s') - E(s))/T)$ at temperature $T$. High temperatures encourage exploration across energy barriers; as the temperature decreases, the search increasingly favors improvements. The best state is tracked separately because the current state may later move to a worse one.

\begin{compactitem}
  \item \inlinecode{anneal\_min(initial, energy, rand\_neighbor, rng, ...)} returns \inlinecode{(best\_energy, best\_state)}. The callable \inlinecode{energy(state)} returns the state's energy, while \inlinecode{rand\_neighbor(state, temp, rng)} returns a randomly chosen nearby state. The optional temperature arguments default to \inlinecode{temp\_start = 1000}, \inlinecode{temp\_end = 1e-6}, and \inlinecode{cooling\_rate = 0.995}. \inlinecode{temp\_end} must be positive and less than \inlinecode{temp\_start}, while \inlinecode{cooling\_rate} must be strictly between 0 and 1.
\end{compactitem}

The geometric schedule repeatedly multiplies the temperature by \inlinecode{cooling\_rate}, performing $k = \lceil \log(T_{end}/T_{start}) / \log(r) \rceil$ iterations. Temperature has the same scale as energy: multiplying every energy by a constant requires multiplying both temperature bounds by the same constant to preserve acceptance probabilities. For continuous states, temperature can also control the neighbor step size. Slower cooling, independent restarts, and a time-based stopping condition often improve results, but simulated annealing never proves optimality.

This generic interface copies each proposed state and recomputes its energy. For large permutation states, adapt the loop to mutate and revert moves in place and update the energy from the affected terms instead.

\textbf{Time Complexity}: $O(k(C_E + C_N))$, where $k$ is the number of temperature levels and $C_E$ and $C_N$ are the costs of evaluating the energy and generating a neighboring state.

\textbf{Space Complexity}: $O(S)$ for the current, next, and best states, where $S$ is the size of one state.

\begin{lstlisting}[language=C++]
#include <cassert>
#include <cmath>
#include <random>
#include <utility>

template<typename State, typename Energy, typename Neighbor>
std::pair<double, State> anneal_min(
    State initial, Energy energy, Neighbor rand_neighbor, std::mt19937 &rng,
    const double temp_start = 1000, const double temp_end = 1e-6, const double cooling_rate = 0.995
) {
  assert(0 < temp_end && temp_end < temp_start);
  assert(0 < cooling_rate && cooling_rate < 1);
  std::uniform_real_distribution<double> unit(0.0, 1.0);
  State current = initial, best = initial;
  double current_energy = energy(current), best_energy = current_energy;
  for (double temp = temp_start; temp > temp_end; temp *= cooling_rate) {
    State next = rand_neighbor(current, temp, rng);
    double next_energy = energy(next);
    if (next_energy <= current_energy ||
        unit(rng) < std::exp((current_energy - next_energy) / temp)) {
      current = next;
      current_energy = next_energy;
      if (current_energy < best_energy) {
        best = current;
        best_energy = current_energy;
      }
    }
  }
  return {best_energy, best};
}
\end{lstlisting}
\Needspace*{4\baselineskip}
\textbf{Example Usage}\par
\begin{lstlisting}[language=C++]
#include <algorithm>
#include <vector>
using namespace std;

int main() {
  mt19937 rng(1234567);  // Set your own seed, or use random_device{}().

  // Escape the local minimum near x = 1 and find the lower minimum near x = -1.
  auto energy_1d = [](double x) { return (x * x - 1) * (x * x - 1) + 0.2 * x; };
  auto neighbor_1d = [](double x, double temp, mt19937 &gen) {
    uniform_real_distribution<double> move(-1.0, 1.0);
    return x + move(gen) * sqrt(temp);
  };
  auto [best_energy, best_x] = anneal_min(1.0, energy_1d, neighbor_1d, rng, 2.0, 1e-8);
  assert(best_energy < -0.2 && best_x < -1);

  // Approximate TSP with a permutation state and random 2-opt reversals.
  vector<pair<double, double>> points{{0, 0}, {1, 0}, {2, 0}, {2, 1},
                                      {2, 2}, {1, 2}, {0, 2}, {0, 1}};
  auto tour_length = [&](const vector<int> &tour) {
    double length = 0;
    for (int i = 0; i < static_cast<int>(tour.size()); i++) {
      auto [x1, y1] = points[tour[i]];
      auto [x2, y2] = points[tour[(i + 1) % tour.size()]];
      length += hypot(x1 - x2, y1 - y2);
    }
    return length;
  };
  auto reverse_segment = [](const vector<int> &tour, double, mt19937 &gen) {
    vector<int> next = tour;
    uniform_int_distribution<int> pick(0, static_cast<int>(tour.size()) - 1);
    int lo = pick(gen), hi = pick(gen);
    if (lo > hi) {
      swap(lo, hi);
    }
    reverse(next.begin() + lo, next.begin() + hi + 1);
    return next;
  };
  vector<int> initial_tour{0, 4, 2, 6, 1, 5, 3, 7};
  auto [best_length, best_tour] =
      anneal_min(initial_tour, tour_length, reverse_segment, rng, 5.0, 1e-6, 0.995);
  assert(fabs(best_length - 8) < 1e-9);
  assert(fabs(tour_length(best_tour) - best_length) < 1e-9);
  return 0;
}
\end{lstlisting}

\section{\texorpdfstring{Dynamic Programming Tricks}{5.3 Dynamic Programming Tricks}}
\setcounter{section}{3}
\setcounter{subsection}{0}
\subsection{Divide-and-Conquer DP Optimization}
Computes one layer of a dynamic program of the form \inlinecode{dp\_cur[i] = min(dp\_prev[k] + cost(k, i))}, where $0 \leq {\inlinecodemath{k}} \leq {\inlinecodemath{i}}$. Let \inlinecode{opt[i]} be the smallest index \inlinecode{k} attaining this minimum. Divide-and-conquer optimization applies when \inlinecode{opt[i]} $\leq$ \inlinecode{opt[i + 1]}, so the best transition index moves only to the right as \inlinecode{i} increases.

To compute \inlinecode{dp\_cur[l..r]}, evaluate its midpoint and record its best transition as \inlinecode{best\_k}. Monotonicity restricts every optimum in the left half to indices at most \inlinecode{best\_k}, and every optimum in the right half to indices at least \inlinecode{best\_k}. Recursively applying these bounds avoids checking every transition for every state. The caller must verify the required monotonicity property.

\begin{compactitem}
  \item \inlinecode{compute\_dp\_layer(dp\_prev, dp\_cur, l, r, opt\_l, opt\_r, cost)} fills \inlinecode{dp\_cur[l..r]} using candidate transition indices in [\inlinecode{opt\_l}, \inlinecode{opt\_r}]. The template parameter \inlinecode{cost} must be callable such that \inlinecode{cost(k, i)} returns the transition cost from previous state \inlinecode{k} to current state \inlinecode{i}.
\end{compactitem}

\textbf{Time Complexity}: $O(n \log n)$ calls to \inlinecode{cost(k, i)} per layer when each state has $O(n)$ possible transitions.

\textbf{Space Complexity}: $O(\log n)$ auxiliary stack space.

\begin{lstlisting}[language=C++]
#include <algorithm>
#include <cstdint>
#include <vector>

const int64_t INF = (1LL << 62);

template<typename Cost>
void compute_dp_layer(
    const std::vector<int64_t> &dp_prev, std::vector<int64_t> &dp_cur, int l, int r, int opt_l,
    int opt_r, Cost cost
) {
  if (l > r) {
    return;
  }
  int mid = l + (r - l) / 2;
  int64_t best = INF;
  int best_k = opt_l;
  int upper = std::min(mid, opt_r);
  for (int k = opt_l; k <= upper; k++) {
    int64_t candidate = dp_prev[k] + cost(k, mid);  // Overflow warning!
    if (candidate < best) {
      best = candidate;
      best_k = k;
    }
  }
  dp_cur[mid] = best;
  compute_dp_layer(dp_prev, dp_cur, l, mid - 1, opt_l, best_k, cost);
  compute_dp_layer(dp_prev, dp_cur, mid + 1, r, best_k, opt_r, cost);
}
\end{lstlisting}
\Needspace*{4\baselineskip}
\textbf{Example Usage}\par
\begin{lstlisting}[language=C++]
#include <cassert>
using namespace std;

struct SquareSegmentCost {
  vector<int64_t> prefix;

  explicit SquareSegmentCost(const vector<int> &a) : prefix(a.size() + 1) {
    for (int i = 0; i < static_cast<int>(a.size()); i++) {
      prefix[i + 1] = prefix[i] + a[i];
    }
  }

  int64_t operator()(int k, int i) const {
    // Cost of placing the half-open segment a[k..i) in one group.
    int64_t sum = prefix[i] - prefix[k];
    return sum * sum;
  }
};

int main() {
  vector<int> a{1, 2, 3, 4};
  int n = static_cast<int>(a.size());
  SquareSegmentCost cost(a);

  // Base layer for zero groups: only the empty prefix is reachable.
  vector<int64_t> dp_prev(n + 1, INF), dp_cur(n + 1, INF);
  dp_prev[0] = 0;

  // The first layer places each prefix a[0..j) into one group.
  compute_dp_layer(dp_prev, dp_cur, 1, n, 0, n - 1, cost);
  assert(dp_cur[4] == 100);  // One group: (1 + 2 + 3 + 4)^2.

  // The second layer chooses the best split between two groups.
  vector<int64_t> dp_two(n + 1, INF);
  compute_dp_layer(dp_cur, dp_two, 1, n, 0, n - 1, cost);
  assert(dp_two[4] == 52);  // (1 + 2 + 3)^2 + 4^2.
  return 0;
}
\end{lstlisting}
\subsection{SMAWK Row Minima}
Finds the position of the minimum in every row of a totally monotone matrix in time linear in its dimensions, without ever materializing the matrix. A matrix is totally monotone (for row minima) when the column index of each row's leftmost minimum never decreases as the row index increases, and this holds not just for the whole matrix but for every submatrix. Monge matrices, where $M_{i,j} + M_{i+1,j+1} \leq M_{i,j+1} + M_{i+1,j}$, are the most common source of such matrices and arise throughout dynamic programming on intervals and partitions.

The SMAWK algorithm interleaves two steps. REDUCE discards columns that cannot hold the minimum of any remaining row, leaving at most as many candidate columns as active rows. INTERPOLATE recurses on the odd-indexed active rows, then recovers each even-indexed active row's minimum by scanning only the candidate-column range bounded by its already-solved neighbors; total monotonicity guarantees those scans advance monotonically and cost $O(R + C)$ overall, where $R$ is the number of rows and $C$ is the number of columns in the original matrix. This removes the logarithmic factor of divide-and-conquer DP optimization, but unlike that technique it is purely offline: every matrix entry must be available on demand from the oracle, so it cannot be used when an entry depends on a row minimum computed earlier in the same pass.

\begin{compactitem}
  \item \inlinecode{smawk\_row\_minima(rows, cols, get)} returns a vector \inlinecode{arg} of length \inlinecode{rows}, where \inlinecode{arg[r]} is the column index minimizing \inlinecode{get(r, c)} over $0 \leq {\inlinecodemath{c}} < {\inlinecodemath{cols}}$ (on ties, the smallest qualifying column is chosen). The input matrix is supplied as a totally monotone function \inlinecode{get(r, c)} with bounds \inlinecode{rows} $\geq 0$ and \inlinecode{cols} $\geq 1$.
\end{compactitem}

\textbf{Time Complexity}: $O(R + C)$ evaluations of \inlinecode{get} per call, where $R$ is \inlinecode{rows} and $C$ is \inlinecode{cols}.

\textbf{Space Complexity}: $O(R + C)$ auxiliary.

\begin{lstlisting}[language=C++]
#include <numeric>
#include <vector>

template<typename Get>
std::vector<int> smawk_rec(
    const Get &get, const std::vector<int> &active_rows, const std::vector<int> &candidate_cols
) {
  int num_rows = static_cast<int>(active_rows.size());
  std::vector<int> arg(num_rows, -1);
  if (num_rows == 0) {
    return arg;
  }
  // Base case: a single row or column is cheap to scan directly.
  if (num_rows <= 2 || candidate_cols.size() <= 2) {
    for (int row_pos = 0; row_pos < num_rows; row_pos++) {
      for (int c : candidate_cols) {
        if (arg[row_pos] == -1 ||
            get(active_rows[row_pos], c) < get(active_rows[row_pos], arg[row_pos])) {
          arg[row_pos] = c;
        }
      }
    }
    return arg;
  }
  // REDUCE: keep only columns that can be the minimum of some row.
  std::vector<int> kept;
  if (static_cast<int>(candidate_cols.size()) > num_rows) {
    for (int c : candidate_cols) {
      while (!kept.empty()) {
        int r = active_rows[static_cast<int>(kept.size()) - 1];
        if (get(r, kept.back()) <= get(r, c)) {
          break;
        }
        kept.pop_back();
      }
      if (static_cast<int>(kept.size()) < num_rows) {
        kept.push_back(c);
      }
    }
  } else {
    kept = candidate_cols;
  }
  // INTERPOLATE: solve the odd rows recursively, then fill the even rows between known bounds.
  std::vector<int> odd_rows;
  for (int row_pos = 1; row_pos < num_rows; row_pos += 2) {
    odd_rows.push_back(active_rows[row_pos]);
  }
  std::vector<int> odd_arg = smawk_rec(get, odd_rows, kept);
  for (int odd_pos = 0; odd_pos < static_cast<int>(odd_arg.size()); odd_pos++) {
    arg[2 * odd_pos + 1] = odd_arg[odd_pos];
  }
  int col_pos = 0;
  for (int row_pos = 0; row_pos < num_rows; row_pos += 2) {
    int hi = (row_pos + 1 < num_rows) ? arg[row_pos + 1] : kept.back();
    int r = active_rows[row_pos];
    while (true) {
      int c = kept[col_pos];
      if (arg[row_pos] == -1 || get(r, c) < get(r, arg[row_pos])) {
        arg[row_pos] = c;
      }
      if (c == hi) {
        break;  // Stop at the upper bound; the next even row starts scanning from here.
      }
      col_pos++;
    }
  }
  return arg;
}

template<typename Get>
std::vector<int> smawk_row_minima(int rows, int cols, const Get &get) {
  std::vector<int> row_ids(rows), col_ids(cols);
  std::iota(row_ids.begin(), row_ids.end(), 0);
  std::iota(col_ids.begin(), col_ids.end(), 0);
  return smawk_rec(get, row_ids, col_ids);
}
\end{lstlisting}
\Needspace*{4\baselineskip}
\textbf{Example Usage}\par
\begin{lstlisting}[language=C++]
#include <cassert>
#include <vector>
using namespace std;

int main() {
  // An explicit totally monotone matrix; row minima sit at non-decreasing columns.
  vector<vector<int>> m{
      {10, 17, 13, 28, 23},
      {17, 22, 16, 29, 23},
      {24, 28, 22, 34, 24},
      {11, 13,  6, 17,  7},
      {45, 44, 32, 37, 23},
  };
  auto get = [&](int r, int c) { return m[r][c]; };
  vector<int> arg = smawk_row_minima(5, 5, get);

  // Cross-check each reported column against a brute-force row scan.
  for (int r = 0; r < 5; r++) {
    int best = 0;
    for (int c = 1; c < 5; c++) {
      if (m[r][c] < m[r][best]) {
        best = c;
      }
    }
    assert(m[r][arg[r]] == m[r][best]);  // Same minimum value (smallest index on ties).
  }
  return 0;
}
\end{lstlisting}
\subsection{Knuth DP Optimization}
Computes minimum costs for every half-open interval $[{\inlinecodemath{l}}, {\inlinecodemath{r}})$ using the recurrence \inlinecode{dp[l][r] = min(dp[l][k] + dp[k][r]) + cost(l, r)}. Each state chooses a split \inlinecode{k} strictly inside the interval and pays \inlinecode{cost(l, r)} after solving the two resulting subintervals.

A direct implementation checks every split and takes $O(n^{3})$. Knuth optimization applies when the smallest optimal splits satisfy \inlinecode{opt[l][r - 1]} $\leq$ \inlinecode{opt[l][r]} $\leq$ \inlinecode{opt[l + 1][r]}. The two neighboring states therefore bound the splits that must be checked for \inlinecode{dp[l][r]}, reducing the total number of candidate checks to $O(n^{2})$. Common applications include optimal binary search trees and some range merging problems. The caller must verify the required monotonicity property for the chosen \inlinecode{cost(l, r)}; the quadrangle inequality and interval monotonicity of the cost are common sufficient conditions.

\begin{compactitem}
  \item \inlinecode{knuth\_interval\_dp(n, cost, \&opt)} computes minimum costs for all half-open intervals $[{\inlinecodemath{l}}, {\inlinecodemath{r}})$ over \inlinecode{n} items. The template parameter \inlinecode{cost} must be callable such that \inlinecode{cost(l, r)} returns the interval cost added after choosing the best split. If the optional pointer \inlinecode{opt} is supplied, it is filled with the chosen split points.
\end{compactitem}

\textbf{Time Complexity}: $O(n^{2})$ calls to \inlinecode{cost(l, r)} and $O(n^{2})$ candidate split checks.

\textbf{Space Complexity}: $O(n^{2})$ for the \inlinecode{dp} and \inlinecode{opt} tables.

\begin{lstlisting}[language=C++]
#include <algorithm>
#include <cstdint>
#include <vector>

const int64_t INF = (1LL << 62);

template<typename Cost>
std::vector<std::vector<int64_t>> knuth_interval_dp(
    int n, Cost cost, std::vector<std::vector<int>> *out = nullptr
) {
  std::vector<std::vector<int64_t>> dp(n + 1, std::vector<int64_t>(n + 1, 0));
  std::vector<std::vector<int>> opt(n + 1, std::vector<int>(n + 1, 0));
  // Empty and one-item intervals have cost 0.
  for (int i = 0; i <= n; i++) {
    opt[i][i] = i;
    if (i < n) {
      opt[i][i + 1] = i + 1;
    }
  }
  for (int len = 2; len <= n; len++) {
    for (int l = 0; l + len <= n; l++) {
      int r = l + len;
      dp[l][r] = INF;
      int start = std::max(opt[l][r - 1], l + 1);
      int finish = std::min(opt[l + 1][r], r - 1);
      int64_t interval_cost = cost(l, r);
      for (int k = start; k <= finish; k++) {
        int64_t candidate = dp[l][k] + dp[k][r] + interval_cost;  // Overflow warning!
        if (candidate < dp[l][r]) {
          dp[l][r] = candidate;
          opt[l][r] = k;
        }
      }
    }
  }
  if (out != nullptr) {
    *out = opt;
  }
  return dp;
}
\end{lstlisting}
\Needspace*{4\baselineskip}
\textbf{Example Usage}\par
\begin{lstlisting}[language=C++]
#include <cassert>
using namespace std;

struct MergeCost {
  vector<int64_t> prefix;

  explicit MergeCost(const vector<int> &a) : prefix(a.size() + 1) {
    for (int i = 0; i < static_cast<int>(a.size()); i++) {
      prefix[i + 1] = prefix[i] + a[i];
    }
  }

  int64_t operator()(int l, int r) const {
    // Cost of merging the half-open interval a[l..r) into one group.
    return prefix[r] - prefix[l];
  }
};

int main() {
  vector<int> a{1, 2, 3, 4};
  vector<vector<int>> opt;
  vector<vector<int64_t>> dp = knuth_interval_dp(a.size(), MergeCost(a), &opt);
  assert(dp[0][4] == 19);  // Merge 1 + 2, then 3 + 3, then 6 + 4.
  assert(opt[0][4] == 3);  // The final merge splits [1, 2, 3] from [4].
  return 0;
}
\end{lstlisting}
\subsection{Lagrangian Relaxation (Aliens Trick)}
Given an optimization problem with a constraint on the number of chosen objects, Lagrangian relaxation moves that count into the objective by charging a penalty per object. If the relaxed problem can be solved quickly and the number chosen is monotone as the penalty changes, binary search recovers the best value for the desired count.

For maximization, solve the relaxed problem as \inlinecode{score - penalty*count}, breaking ties toward larger \inlinecode{count}. As \inlinecode{penalty} increases, the optimal count never increases. The largest penalty whose relaxed optimum still chooses at least \inlinecode{target\_count} objects gives the exact constrained score after adding back \inlinecode{penalty*target\_count}.

\begin{compactitem}
  \item \inlinecode{lagrangian\_maximize(target\_count, lo, hi, solve)} returns the maximum original score using exactly \inlinecode{target\_count} objects. The callable \inlinecode{solve(penalty)} must return \inlinecode{relaxed\_score} and \inlinecode{count} as a pair and must tie-break toward larger \inlinecode{count}. Penalties are searched over the half-open range $[{\inlinecodemath{lo}}, {\inlinecodemath{hi}})$; \inlinecode{lo} must choose at least \inlinecode{target\_count} objects, while \inlinecode{hi} must be greater than the last penalty that does so.
\end{compactitem}

This technique is often called the Aliens trick. Widen the penalty range when in doubt.

\textbf{Time Complexity}: $O(\log n)$ calls to \inlinecode{solve()} per call, where $n$ is the distance between \inlinecode{lo} and \inlinecode{hi}.

\textbf{Space Complexity}: $O(1)$ auxiliary.

\begin{lstlisting}[language=C++]
#include <cassert>
#include <cstdint>
#include <utility>
#include <vector>

template<typename Solve>
int64_t lagrangian_maximize(int target_count, int64_t lo, int64_t hi, Solve solve) {
  // The chosen count is nonincreasing, so find the last penalty that still chooses enough objects.
  while (lo < hi) {
    int64_t mid = lo + (hi - lo) / 2;
    if (solve(mid).second >= target_count) {
      lo = mid + 1;
    } else {
      hi = mid;
    }
  }
  int64_t penalty = lo - 1;
  auto [relaxed_score, count] = solve(penalty);
  assert(count >= target_count);
  return relaxed_score + penalty * target_count;  // Overflow warning!
}
\end{lstlisting}
\Needspace*{4\baselineskip}
\textbf{Example Usage}\par
\begin{lstlisting}[language=C++]
int main() {
  std::vector<int64_t> value{10, 7, 3, 2};
  auto choose_profitable = [&](int64_t penalty) {
    int64_t score = 0;
    int count = 0;
    for (int64_t x : value) {
      int64_t gain = x - penalty;
      if (gain >= 0) {  // Include zero gains to break ties toward larger counts.
        score += gain;
        count++;
      }
    }
    return std::pair<int64_t, int>{score, count};
  };

  assert(lagrangian_maximize(2, 0, 20, choose_profitable) == 17);
  assert(lagrangian_maximize(3, 0, 20, choose_profitable) == 20);
  return 0;
}
\end{lstlisting}
\subsection{Convex Hull Trick (Semi-Dynamic)}
Given a set of pairs $(m, b)$ specifying lines of the form $y = mx + b$, answer queries at specified $x$-coordinates, each asking for the minimum $y$-value over all given lines. This is useful for dynamic programming recurrences of the form \inlinecode{dp[i] = min(m[j] * x[i] + b[j])}. Only the lower envelope of the lines can ever answer a query, so each added line pops previously stored lines that it renders useless, and ascending queries advance a pointer along the envelope so every line is visited at most twice.

The following implementation is a concise, semi-dynamic version of the convex hull optimization technique. It supports an interlaced sequence of \inlinecode{add\_line()} and \inlinecode{query()} calls, as long as the preconditions of descending \inlinecode{m} and ascending \inlinecode{x} are satisfied. As a result, it may be necessary to sort the lines and queries before calling the functions. In that case, the overall time complexity will be dominated by the sorting step.

\begin{compactitem}
  \item \inlinecode{SemiDynamicCHT()} constructs an empty hull.
  \item \inlinecode{add\_line(m, b)} inserts the line $y = mx + b$. The caller must ensure that slope \inlinecode{m} is less than or equal to the slope of every line added so far.
  \item \inlinecode{query(x)} returns the minimum $y$-value among all inserted lines at coordinate \inlinecode{x}. The caller must ensure that query coordinates are nondecreasing across calls. At least one line must have been inserted.
\end{compactitem}

Overflow warning: \inlinecode{add\_line()} compares intersections by cross-multiplying slope and intercept differences, a product on the order of the squared coefficient magnitude. For large \inlinecode{m}/\inlinecode{b} (roughly beyond $10^9$ with \inlinecode{int64\_t}), cast that comparison to \inlinecode{\_\_int128}. \inlinecode{query()} evaluates \inlinecode{m * x + b}, which also needs to fit in \inlinecode{int64\_t}.

\textbf{Time Complexity}: $O(n + q)$ for any interlaced sequence containing $n$ calls to \inlinecode{add\_line()} and $q$ calls to \inlinecode{query()}. Each line is removed or passed by the query pointer at most once, so both operations take amortized $O(1)$ time.

\Needspace*{3\baselineskip}
\textbf{Space Complexity}:
\begin{compactitem}
  \item $O(n)$ for storage of the lines.
  \item $O(1)$ auxiliary for \inlinecode{add\_line()} and \inlinecode{query()}.
\end{compactitem}

\begin{lstlisting}[language=C++]
#include <cassert>
#include <cstdint>
#include <vector>

class SemiDynamicCHT {
  std::vector<int64_t> slopes, intercepts;
  int ptr = 0;

 public:
  void add_line(int64_t m, int64_t b) {
    int len = static_cast<int>(slopes.size());
    if (len > 0 && slopes.back() == m) {
      if (intercepts.back() <= b) {
        return;
      }
      slopes.pop_back();
      intercepts.pop_back();
      len--;
      if (ptr > len) {
        ptr = len;
      }
    }
    // Overflow warning!
    while (len > 1 && (intercepts[len - 2] - intercepts[len - 1]) * (m - slopes[len - 1]) >=
                          (intercepts[len - 1] - b) * (slopes[len - 1] - slopes[len - 2])) {
      len--;
    }
    slopes.resize(len);
    intercepts.resize(len);
    slopes.push_back(m);
    intercepts.push_back(b);
  }

  int64_t query(int64_t x) {
    assert(!slopes.empty());
    if (ptr >= static_cast<int>(slopes.size())) {
      ptr = static_cast<int>(slopes.size()) - 1;
    }
    // Overflow warning!
    while (ptr + 1 < static_cast<int>(slopes.size()) &&
           slopes[ptr + 1] * x + intercepts[ptr + 1] <= slopes[ptr] * x + intercepts[ptr]) {
      ptr++;
    }
    return slopes[ptr] * x + intercepts[ptr];
  }
};
\end{lstlisting}
\Needspace*{4\baselineskip}
\textbf{Example Usage}\par
\begin{lstlisting}[language=C++]
#include <cassert>

int main() {
  SemiDynamicCHT cht;
  cht.add_line(3, 0);
  cht.add_line(2, 1);
  cht.add_line(1, 2);
  cht.add_line(0, 6);
  assert(cht.query(0) == 0);
  assert(cht.query(1) == 3);
  assert(cht.query(2) == 4);
  assert(cht.query(3) == 5);
  cht.add_line(0, 7);  // Same slope, worse intercept; ignored.
  assert(cht.query(4) == 6);
  cht.add_line(0, 5);  // Same slope, better intercept; replaces the old constant line.
  assert(cht.query(5) == 5);
  return 0;
}
\end{lstlisting}
\subsection{Convex Hull Trick (Fully-Dynamic)}
Maintains a dynamic set of lines $y = mx + b$ and answers minimum value queries at integer points. Lines and queries may arrive in arbitrary order, making this useful for dynamic programming recurrences of the form \inlinecode{dp[i] = min(m[j] * x[i] + b[j])} without monotone slopes or query coordinates. The same interface can answer maximum queries when constructed with \inlinecode{query\_max = true}.

This fully dynamic convex hull stores the envelope in a self-balancing binary search tree (\inlinecode{std::set}). Inserting a line removes any neighbors it dominates, and each remaining line records the last coordinate where it is optimal, so a query is one tree lookup. Unlike a Li Chao tree, this structure requires no fixed coordinate domain and its complexity depends on the number of lines rather than the domain width. Prefer it when the domain is unknown or very large, or when maximum queries are needed; a Li Chao tree is often simpler when a manageable domain is known.

\begin{compactitem}
  \item \inlinecode{HullOptimizer(query\_max = false)} constructs an empty hull. By default, \inlinecode{query(x)} minimizes; if \inlinecode{query\_max} is true, \inlinecode{query(x)} maximizes.
  \item \inlinecode{add\_line(m, b)} inserts line $y = mx + b$. Lines may be added in any order.
  \item \inlinecode{query(x)} returns the best $y$-value among all inserted lines at coordinate \inlinecode{x}. At least one line must have been inserted, and query coordinates may be supplied in any order.
\end{compactitem}

Overflow warning: minimization negates the coefficients, border updates subtract slopes/intercepts, and \inlinecode{query()} evaluates \inlinecode{m * x + b}, so those intermediate values must fit in \inlinecode{int64\_t}.

\textbf{Time Complexity}: $O(\log n)$ amortized per call to \inlinecode{add\_line()} and $O(\log n)$ per call to \inlinecode{query()}, where $n$ is the number of lines added.

\Needspace*{3\baselineskip}
\textbf{Space Complexity}:
\begin{compactitem}
  \item $O(n)$ for storage of the lines.
  \item $O(1)$ auxiliary for \inlinecode{add\_line()} and \inlinecode{query()}.
\end{compactitem}

\begin{lstlisting}[language=C++]
#include <cassert>
#include <cstdint>
#include <set>

class HullOptimizer {
  struct Line {
    int64_t m, b;
    mutable int64_t xhi;
    bool is_query;

    Line(int64_t m, int64_t b, int64_t xhi = 0, bool is_query = false)
        : m(m), b(b), xhi(xhi), is_query(is_query) {}

    bool operator<(const Line &l) const { return l.is_query ? xhi < l.xhi : m < l.m; }
  };

  std::multiset<Line> hull;
  bool query_max;

  using hulliter = std::multiset<Line>::iterator;

  bool update_border(hulliter x, hulliter y) {
    if (y == hull.end()) {
      x->xhi = INT64_MAX;
      return false;
    }
    if (x->m == y->m) {
      x->xhi = (x->b > y->b) ? INT64_MAX : INT64_MIN;
    } else {
      int64_t a = y->b - x->b, b = x->m - y->m;  // Overflow warning!
      x->xhi = a / b - ((a ^ b) < 0 && a % b);
    }
    return x->xhi >= y->xhi;
  }

 public:
  explicit HullOptimizer(bool query_max = false) : query_max(query_max) {}

  void add_line(int64_t m, int64_t b) {
    if (!query_max) {
      m = -m;
      b = -b;
    }
    auto z = hull.insert(Line(m, b));
    auto y = z++;
    auto x = y;
    while (update_border(y, z)) {
      z = hull.erase(z);
    }
    if (x != hull.begin() && update_border(--x, y)) {
      update_border(x, y = hull.erase(y));
    }
    while ((y = x) != hull.begin() && (--x)->xhi >= y->xhi) {
      update_border(x, hull.erase(y));
    }
  }

  int64_t query(int64_t x) const {
    assert(!hull.empty());
    Line q(0, 0, x, true);
    auto it = hull.lower_bound(q);
    int64_t res = it->m * x + it->b;  // Overflow warning!
    return query_max ? res : -res;
  }
};
\end{lstlisting}
\Needspace*{4\baselineskip}
\textbf{Example Usage}\par
\begin{lstlisting}[language=C++]
#include <cassert>

int main() {
  HullOptimizer h;
  h.add_line(3, 0);
  h.add_line(0, 6);
  h.add_line(1, 2);
  h.add_line(2, 1);
  // Minimize among y = 3x, 6, x + 2, and 2x + 1.
  assert(h.query(0) == 0);
  assert(h.query(2) == 4);
  assert(h.query(1) == 3);
  assert(h.query(3) == 5);

  HullOptimizer mx(true);
  mx.add_line(3, 0);
  mx.add_line(0, 6);
  mx.add_line(1, 2);
  // Same interface can maximize when constructed with query_max = true.
  assert(mx.query(0) == 6);
  assert(mx.query(3) == 9);
  return 0;
}
\end{lstlisting}
\subsection{Li Chao Tree}
Maintains a dynamic set of lines $y = mx + b$ and answers minimum value queries at integer points. Lines and queries may arrive in arbitrary order, making this useful for dynamic programming recurrences of the form \inlinecode{dp[i] = min(m[j] * x[i] + b[j])} without monotone slopes or query coordinates.

A Li Chao tree fixes an inclusive coordinate domain $[{\inlinecodemath{lo}}, {\inlinecodemath{hi}}]$ and stores one line at each node of a dynamic segment tree. Inserting a line keeps the line that wins at the midpoint and recurses only into the half where the other line may still win; a query takes the best value along one root-to-leaf path. Prefer this structure when the domain is known and manageable. The fully dynamic convex hull in the previous section needs no fixed domain and also supports maximum queries, while a Li Chao tree has a simpler invariant and is easier to extend to lines active only on subintervals.

\begin{compactitem}
  \item \inlinecode{LiChaoTree(lo, hi)} constructs an empty tree over integer domain $[{\inlinecodemath{lo}}, {\inlinecodemath{hi}}]$.
  \item \inlinecode{add\_line(m, b)} inserts line $y = mx + b$. Lines may be added in any order.
  \item \inlinecode{query(x)} returns the minimum $y$-value among all inserted lines at coordinate \inlinecode{x}. At least one line must have been inserted, and query coordinates may be supplied in any order.
\end{compactitem}

Overflow warning: each comparison and query evaluates \inlinecode{m * x + b}, which must fit in \inlinecode{int64\_t}.

\textbf{Time Complexity}: $O(\log d)$ per call to \inlinecode{add\_line(m, b)} and \inlinecode{query(x)}, where $d$ is the distance between \inlinecode{lo} and \inlinecode{hi}.

\Needspace*{3\baselineskip}
\textbf{Space Complexity}:
\begin{compactitem}
  \item $O(n)$ node storage for $n$ lines inserted, since each insertion creates at most one node.
  \item $O(\log d)$ auxiliary stack space per operation.
\end{compactitem}

\begin{lstlisting}[language=C++]
#include <algorithm>
#include <cassert>
#include <cstdint>
#include <utility>

const int64_t INF = INT64_MAX / 4;

class LiChaoTree {
  struct Line {
    int64_t m, b;
    Line(int64_t m = 0, int64_t b = INF) : m(m), b(b) {}
    int64_t eval(int64_t x) const { return m * x + b; }  // Overflow warning!
  };

  struct Node {
    Line f;
    Node *left, *right;
    explicit Node(const Line &f) : f(f), left(nullptr), right(nullptr) {}
  };

  Node *root;
  int64_t lo, hi;

  static void clean_up(Node *n) {
    if (n == nullptr) {
      return;
    }
    clean_up(n->left);
    clean_up(n->right);
    delete n;
  }

  static void add_line(Node *&n, int64_t l, int64_t r, Line f) {
    if (n == nullptr) {
      n = new Node(f);
      return;
    }
    int64_t mid = l + (r - l) / 2;
    bool left_better = f.eval(l) < n->f.eval(l);
    bool mid_better = f.eval(mid) < n->f.eval(mid);
    if (mid_better) {
      std::swap(f, n->f);
    }
    if (l == r) {
      return;
    }
    if (left_better != mid_better) {
      add_line(n->left, l, mid, f);
    } else {
      add_line(n->right, mid + 1, r, f);
    }
  }

  static int64_t query(Node *n, int64_t l, int64_t r, int64_t x) {
    if (n == nullptr) {
      return INF;
    }
    int64_t res = n->f.eval(x);
    if (l == r) {
      return res;
    }
    int64_t mid = l + (r - l) / 2;
    if (x <= mid) {
      return std::min(res, query(n->left, l, mid, x));
    }
    return std::min(res, query(n->right, mid + 1, r, x));
  }

 public:
  LiChaoTree(int64_t lo, int64_t hi) : root(nullptr), lo(lo), hi(hi) {}

  ~LiChaoTree() { clean_up(root); }
  LiChaoTree(const LiChaoTree &) = delete;
  LiChaoTree &operator=(const LiChaoTree &) = delete;
  void add_line(int64_t m, int64_t b) { add_line(root, lo, hi, Line(m, b)); }

  int64_t query(int64_t x) const {
    assert(root != nullptr);
    assert(lo <= x && x <= hi);
    return query(root, lo, hi, x);
  }
};
\end{lstlisting}
\Needspace*{4\baselineskip}
\textbf{Example Usage}\par
\begin{lstlisting}[language=C++]
#include <cassert>

int main() {
  LiChaoTree h(-10, 10);
  h.add_line(3, 0);
  h.add_line(0, 6);
  h.add_line(1, 2);
  h.add_line(2, 1);
  // Minimize among y = 3x, 6, x + 2, and 2x + 1.
  assert(h.query(0) == 0);
  assert(h.query(2) == 4);
  assert(h.query(1) == 3);
  assert(h.query(3) == 5);
  return 0;
}
\end{lstlisting}

\section{\texorpdfstring{Numerical Methods}{5.4 Numerical Methods}}
\setcounter{section}{4}
\setcounter{subsection}{0}
\subsection{Root Finding (Bracketing)}
Finds an $x$ in an interval $[a, b]$ for a continuous function $f$ such that $f(x) = 0$. By the intermediate value theorem, a root must exist in $[a, b]$ if either endpoint is a root or the signs of $f(a)$ and $f(b)$ differ. Each bisection step evaluates the midpoint of the interval and keeps the half whose endpoints still differ in sign, so a root always remains bracketed. After $n$ iterations, the interval width is $(b - a) / 2^n$. Although it is possible to control the error by looping while $b - a$ is greater than an arbitrary epsilon, it is simpler to let the loop run for a desired number of iterations until floating point arithmetic breaks down. 100 iterations is usually sufficient, since the search space will be reduced to $2^{-100}$ (roughly $10^{-30}$) times its original size.

\begin{compactitem}
  \item \inlinecode{bisection\_root(f, a, b, iterations = 100)} returns a root in an interval $[{\inlinecodemath{a}}, {\inlinecodemath{b}}]$ for a continuous function $f$ where either endpoint is a root or the endpoint values have opposite signs, using the bisection method.
  \item \inlinecode{falsi\_illinois\_root(f, a, b, iterations = 100)} returns a root in an interval $[{\inlinecodemath{a}}, {\inlinecodemath{b}}]$ for a continuous function $f$ under the same endpoint conditions, using the Illinois algorithm variant of the false position (a.k.a. regula falsi) method.
\end{compactitem}

\textbf{Time Complexity}: $O(n)$ calls will be made to \inlinecode{f()} in \inlinecode{bisection\_root()} and \inlinecode{falsi\_illinois\_root()}, where $n$ is the number of iterations performed.

\textbf{Space Complexity}: $O(1)$ auxiliary for both operations.

\begin{lstlisting}[language=C++]
#include <cassert>

template<typename Fn>
double bisection_root(Fn f, double a, double b, const int iterations = 100) {
  double fa = f(a), fb = f(b);
  assert(a <= b && ((fa <= 0 && fb >= 0) || (fa >= 0 && fb <= 0)));
  if (fa == 0) {
    return a;
  }
  if (fb == 0) {
    return b;
  }
  double m = a;
  for (int i = 0; i < iterations; i++) {
    m = a + (b - a) / 2;
    double fm = f(m);
    if (fm == 0) {
      return m;
    }
    if ((fa < 0) == (fm < 0)) {
      a = m;
      fa = fm;
    } else {
      b = m;
    }
  }
  return m;
}

template<typename Fn>
double falsi_illinois_root(Fn f, double a, double b, const int iterations = 100) {
  double fa = f(a), fb = f(b);
  assert(a <= b && ((fa <= 0 && fb >= 0) || (fa >= 0 && fb <= 0)));
  if (fa == 0) {
    return a;
  }
  if (fb == 0) {
    return b;
  }
  double m = a;
  int side = 0;
  for (int i = 0; i < iterations; i++) {
    m = (fa * b - fb * a) / (fa - fb);
    double fm = f(m);
    if (fm == 0) {
      break;
    }
    if ((fb < 0) == (fm < 0)) {
      b = m;
      fb = fm;
      if (side < 0) {
        fa /= 2;
      }
      side = -1;
    } else {
      a = m;
      fa = fm;
      if (side > 0) {
        fb /= 2;
      }
      side = 1;
    }
  }
  return m;
}
\end{lstlisting}
\Needspace*{4\baselineskip}
\textbf{Example Usage}\par
\begin{lstlisting}[language=C++]
#include <cmath>

double f(double x) {
  return x * x - 4 * sin(x);
}

int main() {
  assert(fabs(f(bisection_root(f, 1, 3))) < 1e-10);
  assert(fabs(f(falsi_illinois_root(f, 1, 3))) < 1e-10);
  assert(bisection_root(f, 0, 1) == 0);
  assert(falsi_illinois_root(f, 0, 1) == 0);
  return 0;
}
\end{lstlisting}
\subsection{Root Finding (Iteration)}
Finds an $x$ for a continuous function $f$ such that $f(x) = 0$ using iterative approximation by an initial guess that is close to the answer. Each step follows the tangent line at the current guess down to its $x$-intercept, that is, $x \leftarrow x - f(x)/f'(x)$. Newton's method requires an explicit definition of the function's derivative while the secant method starts with two initial guesses and approximates the derivative using the secant slope from the previous iteration. For $n$ iterations and a good initial guess, the methods below compute approximately $2^n$ digits of precision, with the secant method converging approximately $1.6$ times slower than Newton's.

\begin{compactitem}
  \item \inlinecode{newton\_root(f, fprime, x0, eps = 1e-15, iterations = 100)} returns a root $x$ for a function \inlinecode{f} with derivative \inlinecode{fprime} using an initial guess \inlinecode{x0} which should be relatively close to $x$.
  \item \inlinecode{secant\_root(f, x0, x1, eps = 1e-15, iterations = 100)} returns a root $x$ for a function \inlinecode{f} using two initial guesses \inlinecode{x0} and \inlinecode{x1} which should be relatively close to $x$.
\end{compactitem}

\textbf{Time Complexity}: $O(n)$ calls will be made to \inlinecode{f()} in \inlinecode{newton\_root()} and \inlinecode{secant\_root()}, where $n$ is the number of iterations performed.

\textbf{Space Complexity}: $O(1)$ auxiliary for both operations.

\begin{lstlisting}[language=C++]
#include <cmath>
#include <stdexcept>

template<typename Fn, typename Deriv>
double newton_root(
    Fn f, Deriv fprime, double x0, const double eps = 1e-15, const int iterations = 100
) {
  double x = x0, error = eps + 1;
  for (int i = 0; error > eps && i < iterations; i++) {
    double xnew = x - f(x) / fprime(x);
    error = std::fabs(xnew - x);
    x = xnew;
  }
  if (error > eps) {
    throw std::runtime_error("Newton's method failed to converge.");
  }
  return x;
}

template<typename Fn>
double secant_root(
    Fn f, double x0, double x1, const double eps = 1e-15, const int iterations = 100
) {
  double xold = x0, fxold = f(x0), x = x1, error = eps + 1;
  for (int i = 0; error > eps && i < iterations; i++) {
    double fx = f(x);
    double xnew = x - fx * ((x - xold) / (fx - fxold));
    xold = x;
    fxold = fx;
    error = std::fabs(xnew - x);
    x = xnew;
  }
  if (error > eps) {
    throw std::runtime_error("Secant method failed to converge.");
  }
  return x;
}
\end{lstlisting}
\Needspace*{4\baselineskip}
\textbf{Example Usage}\par
\begin{lstlisting}[language=C++]
#include <cassert>

double f(double x) {
  return x * x - 4 * sin(x);
}

double fprime(double x) {
  return 2 * x - 4 * cos(x);
}

int main() {
  assert(fabs(f(newton_root(f, fprime, 3))) < 1e-10);
  assert(fabs(f(secant_root(f, 3, 2))) < 1e-10);
  auto g = [](double x) { return x * x - 2; };
  auto gprime = [](double x) { return 2 * x; };
  assert(fabs(g(newton_root(g, gprime, 1))) < 1e-10);
  return 0;
}
\end{lstlisting}
\subsection{Polynomial Root Finding (Differentiation)}
Finds every real root $x$ of a polynomial $p$ satisfying $p(x) = 0$ by recursively finding the roots of its derivative. Each adjacent pair of local extrema is searched using the bisection method.

The key observation is that derivative roots split the real line into intervals where the polynomial is monotone. A monotone interval contains at most one root, and if an endpoint is a root or the endpoint values have opposite signs, bisection isolates that root. Recursively solving the derivative therefore provides all interval boundaries needed to find every real root in the requested range, including roots of even multiplicity.

\begin{compactitem}
  \item \inlinecode{horner\_eval(p, x)} evaluates the polynomial \inlinecode{p} of degree $d$ (represented as a vector of size $d + 1$ where \inlinecode{p[i]} stores the coefficient for the $x^i$ term) at \inlinecode{x}, using Horner's method.
  \item \inlinecode{find\_one\_root(p, a, b, eps = 1e-15)} returns a root in the interval $[{\inlinecodemath{a}}, {\inlinecodemath{b}}]$ for a polynomial \inlinecode{p} where either endpoint is a root or the endpoint values have opposite signs, using the bisection method. If this precondition is not satisfied, then \inlinecode{NaN} is returned. The root is found to a tolerance of \inlinecode{eps} in absolute or relative error (whichever is reached first).
  \item \inlinecode{find\_all\_roots(p, a = -1e20, b = 1e20, eps = 1e-15)} returns a vector of all roots in the interval $[{\inlinecodemath{a}}, {\inlinecodemath{b}}]$ for a polynomial \inlinecode{p} using the bisection method. The roots are found to a tolerance of \inlinecode{eps} in absolute or relative error (whichever is reached first).
\end{compactitem}

\Needspace*{3\baselineskip}
\textbf{Time Complexity}:
\begin{compactitem}
  \item $O(n)$ per call to \inlinecode{horner\_eval()}, where $n$ is the degree of the polynomial.
  \item $O(nt)$ per call to \inlinecode{find\_one\_root()}, where $n$ is the degree of the polynomial and $t$ is the number of bisection iterations required to reach the requested tolerance.
  \item $O(n^{3} t)$ per call to \inlinecode{find\_all\_roots()}, where $n$ is the degree of the polynomial and $t$ is the number of bisection iterations required to reach the requested tolerance.
\end{compactitem}

\Needspace*{3\baselineskip}
\textbf{Space Complexity}:
\begin{compactitem}
  \item $O(1)$ auxiliary for \inlinecode{horner\_eval()} and \inlinecode{find\_one\_root()}.
  \item $O(n)$ auxiliary stack space and $O(n^{2})$ auxiliary heap space for \inlinecode{find\_all\_roots()}, where $n$ is the degree of the polynomial.
\end{compactitem}

\begin{lstlisting}[language=C++]
#include <cmath>
#include <limits>
#include <vector>

double horner_eval(const std::vector<double> &p, double x) {
  double res = p.back();
  for (int i = static_cast<int>(p.size()) - 2; i >= 0; i--) {
    res = res * x + p[i];
  }
  return res;
}

double find_one_root(const std::vector<double> &p, double a, double b, const double eps = 1e-15) {
  double pa = horner_eval(p, a), pb = horner_eval(p, b);
  if (std::fabs(pa) <= eps) {
    return a;
  }
  if (std::fabs(pb) <= eps) {
    return b;
  }
  bool paneg = pa < 0, pbneg = pb < 0;
  if (paneg == pbneg) {
    return std::numeric_limits<double>::quiet_NaN();
  }
  while (b - a > eps && a * (1 + eps) < b && a < b * (1 + eps)) {
    double m = a + (b - a) / 2;
    if ((horner_eval(p, m) < 0) == paneg) {
      a = m;
    } else {
      b = m;
    }
  }
  return a;
}

std::vector<double> find_all_roots(
    const std::vector<double> &p, double a = -1e20, double b = 1e20, const double eps = 1e-15
) {
  std::vector<double> pprime;
  pprime.reserve(p.size() > 0 ? p.size() - 1 : 0);
  for (int i = 1; i < static_cast<int>(p.size()); i++) {
    pprime.push_back(p[i] * i);
  }
  if (pprime.empty()) {
    return {};
  }
  std::vector<double> res, r = find_all_roots(pprime, a, b, eps);
  r.push_back(b);
  for (int i = 0; i < static_cast<int>(r.size()); i++) {
    double root = find_one_root(p, i == 0 ? a : r[i - 1], r[i], eps);
    if (!std::isnan(root) && (res.empty() || root != res.back())) {
      res.push_back(root);
    }
  }
  return res;
}
\end{lstlisting}
\Needspace*{4\baselineskip}
\textbf{Example Usage}\par
\begin{lstlisting}[language=C++]
#include <cassert>
using namespace std;

int main() {
  {  // -1 + 2x - 6x^2 + 2x^3
    vector<double> p{-1.0, 2.0, -6.0, 2.0};
    vector<double> roots = find_all_roots(p);
    assert(roots.size() == 1 && fabs(horner_eval(p, roots[0])) < 1e-10);
  }
  {  // -20 + 4x + 3x^2
    vector<double> p{-20.0, 4.0, 3.0};
    vector<double> roots = find_all_roots(p);
    assert(roots.size() == 2);
    assert(fabs(horner_eval(p, roots[0])) < 1e-10);
    assert(fabs(horner_eval(p, roots[1])) < 1e-10);
  }
  {  // (x - 1)^2
    vector<double> p{1.0, -2.0, 1.0};
    vector<double> roots = find_all_roots(p);
    assert(roots.size() == 1 && fabs(roots[0] - 1) < 1e-10);
  }
  return 0;
}
\end{lstlisting}
\subsection{Polynomial Root Finding (Laguerre)}
Finds every complex root $x$ for a polynomial $p$ with complex coefficients such that $p(x) = 0$. Laguerre's method is a root-finding iteration with reliable, near-cubic convergence from almost any starting point. Once a root is found it is removed by polynomial deflation, and the iteration repeats on the lower-degree quotient until every root has been extracted.

\begin{compactitem}
  \item \inlinecode{horner\_eval(p, x)} evaluates a complex polynomial $p$ of degree $d$ (represented as a vector of size $d + 1$ where \inlinecode{p[i]} stores the complex coefficient for the $x^i$ term) at \inlinecode{x}, using Horner's method, returning a pair where the first value is the final result $p(x)$ and the second value is the quotient polynomial $q$ from the identity $p(t) = (t - x)q(t) + p(x)$.
  \item \inlinecode{find\_one\_root(p, x0, eps = 1e-15, iterations = 10000)} returns a complex root $x$ for polynomial \inlinecode{p} (represented as a vector of size $d + 1$ where \inlinecode{p[i]} stores the complex coefficient for the $x^i$ term) using an initial guess \inlinecode{x0} which should be relatively close to $x$. The root is found to a tolerance of \inlinecode{eps} in absolute or relative error (whichever is reached first).
  \item \inlinecode{find\_all\_roots(p, eps = 1e-15, iterations = 10000)} returns a vector of all complex roots for a complex polynomial \inlinecode{p}. The roots are found to a tolerance of \inlinecode{eps} in absolute or relative error (whichever is reached first).
\end{compactitem}

\Needspace*{3\baselineskip}
\textbf{Time Complexity}:
\begin{compactitem}
  \item $O(n)$ per call to \inlinecode{horner\_eval()}, where $n$ is the degree of the polynomial.
  \item $O(nt)$ per call to \inlinecode{find\_one\_root()}, where $n$ is the degree of the polynomial and $t$ is the number of iterations performed.
  \item $O(n^{2} t)$ per call to \inlinecode{find\_all\_roots()}, where $n$ is the degree of the polynomial and $t$ is the number of iterations performed per root.
\end{compactitem}

\Needspace*{3\baselineskip}
\textbf{Space Complexity}:
\begin{compactitem}
  \item $O(n)$ auxiliary for \inlinecode{horner\_eval()} and \inlinecode{find\_one\_root()}, where $n$ is the degree of the polynomial.
  \item $O(n)$ auxiliary for \inlinecode{find\_all\_roots()}, where $n$ is the degree of the polynomial.
\end{compactitem}

\begin{lstlisting}[language=C++]
#include <algorithm>
#include <cmath>
#include <complex>
#include <random>
#include <utility>
#include <vector>

using cdouble = std::complex<double>;
using cpoly = std::vector<cdouble>;

std::pair<cdouble, cpoly> horner_eval(const cpoly &p, const cdouble &x) {
  int n = static_cast<int>(p.size());
  cpoly b(std::max(1, n - 1));
  for (int i = n - 1; i > 0; i--) {
    b[i - 1] = p[i] + (i < n - 1 ? b[i] * x : 0);
  }
  return {p[0] + b[0] * x, b};
}

cpoly derivative(const cpoly &p) {
  int n = static_cast<int>(p.size());
  cpoly res(std::max(1, n - 1));
  for (int i = 1; i < n; i++) {
    res[i - 1] = p[i] * cdouble(i);
  }
  return res;
}

int comp(const cdouble &a, const cdouble &b, const double eps = 1e-15) {
  double diff = std::abs(a) - std::abs(b);
  return (diff < -eps) ? -1 : (diff > eps ? 1 : 0);
}

cdouble find_one_root(
    const cpoly &p, const cdouble &x0, const double eps = 1e-15, const int iterations = 10000
) {
  cdouble x = x0;
  int n = static_cast<int>(p.size()) - 1;
  cpoly p1 = derivative(p), p2 = derivative(p1);
  for (int i = 0; i < iterations; i++) {
    cdouble y0 = horner_eval(p, x).first;
    if (comp(y0, 0, eps) == 0) {
      break;
    }
    cdouble g = horner_eval(p1, x).first / y0;
    cdouble h = g * g - horner_eval(p2, x).first / y0;
    cdouble r = std::sqrt(cdouble(n - 1) * (h * cdouble(n) - g * g));
    cdouble d1 = g + r, d2 = g - r;
    cdouble a = cdouble(n) / (comp(d1, d2, eps) > 0 ? d1 : d2);
    x -= a;
    if (comp(a, 0, eps) == 0) {
      break;
    }
  }
  return x;
}

std::vector<cdouble> find_all_roots(
    const cpoly &p, const double eps = 1e-15, const int iterations = 10000
) {
  std::vector<cdouble> res;
  cpoly q = p;
  if (q.size() <= 1) {
    return res;
  }
  static std::mt19937 rng(std::random_device{}());
  std::uniform_real_distribution<double> unit(0.0, 1.0);
  while (q.size() > 2) {
    cdouble z(unit(rng), unit(rng));
    z = find_one_root(p, find_one_root(q, z, eps, iterations), eps, iterations);
    q = horner_eval(q, z).second;
    res.push_back(z);
  }
  res.push_back(-q[0] / q[1]);
  return res;
}
\end{lstlisting}
\Needspace*{4\baselineskip}
\textbf{Example Usage}\par
\begin{lstlisting}[language=C++]
#include <cstdio>
using namespace std;

void print_roots(vector<cdouble> x) {
  sort(x.begin(), x.end(), [](const cdouble &a, const cdouble &b) {
    return abs(a.real() - b.real()) > 0.5e-5 ? a.real() < b.real() : a.imag() < b.imag();
  });
  for (const auto &z : x) {
    double real = abs(z.real()) < 0.5e-5 ? 0 : z.real();
    double imag = abs(z.imag()) < 0.5e-5 ? 0 : z.imag();
    printf("(%.5lf, %.5lf)\n", real, imag);
  }
}

int main() {
  {  // 140 - 13x - 8x^2 + x^3 = (x + 4)(x - 5)(x - 7)
    printf("Roots of 140 - 13x - 8x^2 + x^3:\n");
    cpoly p;
    p.push_back(140);
    p.push_back(-13);
    p.push_back(-8);
    p.push_back(1);
    print_roots(find_all_roots(p));
  }
  {  // (-24+36i) + (-26+12i)x + (-30+40i)x^2 + (-26+12i)x^3 + (-6+4i)x^4
    // = ((2 + 3i)x + 6)(x + i)(2x + (6 + 4i))(xi + 1):
    printf("Roots of ((2 + 3i)x + 6)(x + i)(2x + (6 + 4i))(xi + 1):\n");
    cpoly p;
    p.emplace_back(-24, 36);
    p.emplace_back(-26, 12);
    p.emplace_back(-30, 40);
    p.emplace_back(-26, 12);
    p.emplace_back(-6, 4);
    print_roots(find_all_roots(p));
  }
  return 0;
}
\end{lstlisting}
\begin{exampleoutput}
Roots of 140 - 13x - 8x^2 + x^3:
(-4.00000, 0.00000)
(5.00000, 0.00000)
(7.00000, 0.00000)
Roots of ((2 + 3i)x + 6)(x + i)(2x + (6 + 4i))(xi + 1):
(-3.00000, -2.00000)
(-0.92308, 1.38462)
(0.00000, -1.00000)
(0.00000, 1.00000)
\end{exampleoutput}
\subsection{Polynomial Root Finding (Ehrlich-Aberth)}
Finds every complex root $x$ for a polynomial $p$ such that $p(x) = 0$. The Ehrlich-Aberth method is a simultaneous iteration: every root estimate is updated using both the Newton correction and a repulsion term from the other estimates.

This routine is intended for well-scaled polynomials when all complex roots are wanted. If only real roots are needed, a real-root isolator such as derivative recursion with bisection is usually more reliable. Multiple or tightly clustered roots may converge slowly, and very large or very small root scales may require rescaling the input or using multiprecision arithmetic.

\begin{compactitem}
  \item \inlinecode{eval\_with\_derivative(p, x)} returns a pair $(p(x), p'(x))$ for a polynomial $p$ given as a vector \inlinecode{p} where \inlinecode{p[i]} stores the coefficient for the $x^i$ term.
  \item \inlinecode{find\_all\_roots(p, eps = ROOT\_EPS, iterations = 2000)} returns a vector of all complex roots for a complex polynomial given by the vector of coefficients \inlinecode{p}. A \inlinecode{vector\textless{}dbl\textgreater{}} overload is provided for polynomials with real coefficients. The roots are found to a tolerance of \inlinecode{eps} in absolute or relative error (whichever is reached first), and zero roots are removed exactly before the simultaneous iteration starts.
\end{compactitem}

\Needspace*{3\baselineskip}
\textbf{Time Complexity}:
\begin{compactitem}
  \item $O(n)$ per call to \inlinecode{eval\_with\_derivative(p, x)}, where $n$ is the degree of the polynomial.
  \item $O(n^{2} t)$ per call to \inlinecode{find\_all\_roots(p, eps, iterations)}, where $n$ is the degree of the polynomial and $t$ is the number of iterations required to reach the desired precision.
\end{compactitem}

\Needspace*{3\baselineskip}
\textbf{Space Complexity}:
\begin{compactitem}
  \item $O(1)$ auxiliary for \inlinecode{eval\_with\_derivative(p, x)}.
  \item $O(n)$ auxiliary for \inlinecode{find\_all\_roots(p, eps, iterations)}, where $n$ is the degree of the polynomial.
\end{compactitem}

\begin{lstlisting}[language=C++]
#include <algorithm>
#include <cmath>
#include <complex>
#include <utility>
#include <vector>

using dbl = long double;
using cdbl = std::complex<dbl>;
using cpoly = std::vector<cdbl>;

const dbl PI = acosl(-1.0L);
const dbl ZERO_EPS = 1e-30L;   // Treat coefficients and denominators this small as zero.
const dbl ROOT_EPS = 1e-18L;   // Stop once root updates are this small relative to the root.
const dbl CHECK_EPS = 1e-12L;  // Residual tolerance used by the example assertions.

bool is_zero(const cdbl &z, const dbl eps = ZERO_EPS) {
  return std::abs(z) <= eps;
}

bool is_finite(const cdbl &z) {
  return std::isfinite(z.real()) && std::isfinite(z.imag());
}

std::pair<cdbl, cdbl> eval_with_derivative(const cpoly &p, const cdbl &x) {
  cdbl value = p.back(), derivative = 0;
  for (int i = static_cast<int>(p.size()) - 2; i >= 0; i--) {
    derivative = derivative * x + value;
    value = value * x + p[i];
  }
  return {value, derivative};
}

dbl root_bound(const cpoly &p) {
  dbl res = 0;
  int n = static_cast<int>(p.size()) - 1;
  for (int i = 0; i < n; i++) {
    dbl ratio = std::abs(p[i] / p.back());
    if (ratio > 0) {
      res = std::max(res, powl(ratio, 1.0L / (n - i)));
    }
  }
  return 2 * std::max(static_cast<dbl>(1), res);
}

bool root_less(const cdbl &a, const cdbl &b) {
  if (a.real() != b.real()) {
    return a.real() < b.real();
  }
  return a.imag() < b.imag();
}

cpoly find_all_roots(cpoly p, const dbl eps = ROOT_EPS, const int iterations = 2000) {
  while (!p.empty() && is_zero(p.back())) {
    p.pop_back();
  }
  cpoly roots;
  while (p.size() > 1 && is_zero(p[0])) {
    roots.push_back(0);
    p.erase(p.begin());
  }
  if (p.size() <= 1) {
    return roots;
  }
  dbl scale = 0;
  for (int i = 0; i < static_cast<int>(p.size()); i++) {
    scale = std::max(scale, std::abs(p[i]));
  }
  for (int i = 0; i < static_cast<int>(p.size()); i++) {
    p[i] /= scale;
  }
  int n = static_cast<int>(p.size()) - 1;
  if (n == 1) {
    roots.push_back(-p[0] / p[1]);
    return roots;
  }
  cpoly z(n);
  dbl radius = root_bound(p), offset = PI / (2 * n);
  dbl max_radius = 2 * radius;
  for (int i = 0; i < n; i++) {
    dbl angle = offset + 2 * PI * i / n;
    z[i] = radius * cdbl(cosl(angle), sinl(angle));
  }
  for (int it = 0; it < iterations; it++) {
    bool done = true;
    cpoly next = z;
    for (int i = 0; i < n; i++) {
      auto [fx, dfx] = eval_with_derivative(p, z[i]);
      if (std::abs(fx) <= eps) {
        continue;
      }
      cdbl repulsion = 0;
      for (int j = 0; j < n; j++) {
        if (i != j) {
          cdbl diff = z[i] - z[j];
          if (std::abs(diff) <= eps * eps) {
            dbl angle = 2 * PI * (i + 1) / (n + 1);
            diff = eps * (1 + std::abs(z[i])) * cdbl(cosl(angle), sinl(angle));
          }
          repulsion += cdbl(1) / diff;
        }
      }
      cdbl denom = dfx - fx * repulsion;
      if (is_zero(denom)) {
        done = false;
        continue;
      }
      cdbl step = fx / denom;
      if (!is_finite(step) && !is_zero(dfx)) {
        step = fx / dfx;
      }
      if (!is_finite(step)) {
        done = false;
        continue;
      }
      dbl limit = 2 * max_radius;
      if (std::abs(step) > limit) {
        step *= limit / std::abs(step);
      }
      next[i] = z[i] - step;
      if (!is_finite(next[i])) {
        next[i] = z[i];
        done = false;
        continue;
      }
      if (std::abs(next[i]) > max_radius) {
        next[i] *= max_radius / std::abs(next[i]);
      }
      if (std::abs(step) > eps * (1 + std::abs(next[i]))) {
        done = false;
      }
    }
    z = next;
    if (done) {
      break;
    }
  }
  for (int i = 0; i < n; i++) {
    roots.emplace_back(z[i]);
  }
  std::sort(roots.begin(), roots.end(), root_less);
  return roots;
}

std::vector<cdbl> find_all_roots(
    const std::vector<dbl> &p, const dbl eps = ROOT_EPS, const int iterations = 2000
) {
  cpoly q(p.begin(), p.end());
  return find_all_roots(q, eps, iterations);
}
\end{lstlisting}
\Needspace*{4\baselineskip}
\textbf{Example Usage}\par
\begin{lstlisting}[language=C++]
#include <cassert>
#include <cstdio>
using namespace std;

cdbl eval(const cpoly &p, const cdbl &x) {
  return eval_with_derivative(p, x).first;
}

void print_roots(vector<cdbl> x) {
  sort(x.begin(), x.end(), [](const cdbl &a, const cdbl &b) {
    return abs(a.real() - b.real()) > 0.5e-5 ? a.real() < b.real() : a.imag() < b.imag();
  });
  for (const auto &z : x) {
    dbl real = abs(z.real()) < 0.5e-5 ? 0 : z.real();
    dbl imag = abs(z.imag()) < 0.5e-5 ? 0 : z.imag();
    printf("(%.5Lf, %.5Lf)\n", real, imag);
  }
}

int main() {
  {  // -1 + 2x - 6x^2 + 2x^3
    printf("Roots of -1 + 2x - 6x^2 + 2x^3:\n");
    vector<dbl> p{-1.0, 2.0, -6.0, 2.0};
    vector<cdbl> roots = find_all_roots(p);
    assert(roots.size() == 3);
    for (const auto &root : roots) {
      assert(abs(eval(cpoly(p.begin(), p.end()), root)) < CHECK_EPS);
    }
    print_roots(roots);
  }
  {  // (-24+36i) + (-26+12i)x + (-30+40i)x^2 + (-26+12i)x^3 + (-6+4i)x^4
    // = ((2 + 3i)x + 6)(x + i)(2x + (6 + 4i))(xi + 1):
    printf("Roots of ((2 + 3i)x + 6)(x + i)(2x + (6 + 4i))(xi + 1):\n");
    cpoly p;
    p.emplace_back(-24, 36);
    p.emplace_back(-26, 12);
    p.emplace_back(-30, 40);
    p.emplace_back(-26, 12);
    p.emplace_back(-6, 4);
    vector<cdbl> roots = find_all_roots(p);
    assert(roots.size() == 4);
    for (const auto &root : roots) {
      assert(abs(eval(p, root)) < CHECK_EPS);
    }
    print_roots(roots);
  }
  return 0;
}
\end{lstlisting}
\begin{exampleoutput}
Roots of -1 + 2x - 6x^2 + 2x^3:
(0.15098, -0.40314)
(0.15098, 0.40314)
(2.69805, 0.00000)
Roots of ((2 + 3i)x + 6)(x + i)(2x + (6 + 4i))(xi + 1):
(-3.00000, -2.00000)
(-0.92308, 1.38462)
(0.00000, -1.00000)
(0.00000, 1.00000)
\end{exampleoutput}
\subsection{Integration (Simpson)}
Computes the definite integral from $a$ to $b$ for a continuous function $f$ using the composite Simpson's rule: the interval is split into $n$ equal subintervals and each consecutive pair is approximated by a parabola, giving the approximation $\int_{a}^{b} f(x) \,dx \approx$ $h/3 \cdot [f(x_0) + 4f(x_1) + 2f(x_2) + \ldots + 4f(x_{n-1}) + f(x_n)]$ with step $h = (b - a)/n$.

Unlike adaptive quadrature, the step size here is uniform, so a feature narrower than \inlinecode{h} may be under-resolved. The benefit is a simple, non-recursive routine with predictable cost that is not fooled by the premature-termination heuristic that adaptive Simpson's rule can suffer on functions whose coarse sample points happen to fit a parabola.

\begin{compactitem}
  \item \inlinecode{integrate(f, a, b, n = 1000000)} returns the definite integral of a function \inlinecode{f} from \inlinecode{a} to \inlinecode{b} using \inlinecode{n} subintervals, where \inlinecode{n} must be even. Larger \inlinecode{n} trades runtime for accuracy; the leading error scales with \inlinecode{(b - a) * h\textasciicircum{}4} and the maximum magnitude of the fourth derivative of $f$ over the interval.
\end{compactitem}

\textbf{Time Complexity}: $O(n)$ per call, requiring $n + 1$ evaluations of $f$.

\textbf{Space Complexity}: $O(1)$ auxiliary.

\begin{lstlisting}[language=C++]
#include <cassert>

template<typename Fn>
double integrate(Fn f, double a, double b, int n = 1000000) {
  assert(n > 0 && n % 2 == 0);
  double h = (b - a) / n, s = f(a) + f(b);
  for (int i = 1; i < n; i++) {
    s += f(a + h * i) * (i & 1 ? 4 : 2);
  }
  return s * h / 3;
}
\end{lstlisting}
\Needspace*{4\baselineskip}
\textbf{Example Usage}\par
\begin{lstlisting}[language=C++]
#include <cmath>
using namespace std;

double f(double x) {
  return sin(x);
}

int main() {
  const double PI = acos(-1.0);
  assert(fabs(integrate(f, 0.0, PI / 2) - 1) < 1e-10);
  assert(fabs(integrate(f, PI / 2, 0.0) + 1) < 1e-10);
  assert(fabs(integrate(f, 1.0, 1.0)) < 1e-12);
  return 0;
}
\end{lstlisting}
\subsection{Linear Programming (Simplex)}
Solves a linear programming problem using Dantzig's simplex algorithm. The canonical form of a linear programming problem is to maximize (or minimize) the dot product $cx$, subject to $ax \leq b$ and $x \geq 0$, where $x$ is a vector of unknowns to be solved, $c$ is a vector of coefficients, $a$ is a matrix of linear equation coefficients, and $b$ is a vector of boundary coefficients.

The simplex algorithm walks between vertices of the feasible polytope, at each step pivoting to an adjacent vertex that improves the objective, until no improving move remains (an optimum is reached) or an unbounded improving direction is found.

The implementation uses a two-phase simplex tableau. Phase 1 introduces an artificial variable to find a feasible starting basis, which is essential when some constraint bounds are negative. Phase 2 then optimizes the original objective. The \inlinecode{basis} and \inlinecode{nonbasis} arrays track which variable each tableau row or column currently represents, making solution recovery independent of pivot order.

\begin{compactitem}
  \item \inlinecode{simplex\_solve(a, b, c, \&x, maximize = true, eps = 1e-10)} solves the linear programming problem for an $m$ by $n$ matrix \inlinecode{a} of real values, a length $m$ vector \inlinecode{b}, and a length $n$ vector \inlinecode{c}, returning $0$ if an optimum was found, $-1$ if there are no feasible solutions, or $1$ if the objective is unbounded. Set \inlinecode{maximize = false} to minimize instead, and use \inlinecode{eps} as the pivot and feasibility tolerance. If an optimum is found, \inlinecode{x} is assigned a vector of length \inlinecode{n}, where \inlinecode{x[j]} is the value of the $j$-th variable.
\end{compactitem}

\textbf{Time Complexity}: $O(pmn)$, where $p$ is the total number of pivots across both phases. Each pivot takes $O(mn)$, and $p$ can be exponential in the worst case, although it is usually much smaller in practice.

\textbf{Space Complexity}: $O(m \cdot n)$ auxiliary.

\begin{lstlisting}[language=C++]
#include <cassert>
#include <cmath>
#include <utility>
#include <vector>

int simplex_solve(
    const std::vector<std::vector<double>> &a, const std::vector<double> &b,
    const std::vector<double> &c, std::vector<double> *x, const bool maximize = true,
    const double eps = 1e-10
) {
  int m = static_cast<int>(a.size()), n = static_cast<int>(c.size());
  assert(x != nullptr && b.size() == a.size());
  std::vector<int> basis(m), nonbasis(n + 1);
  std::vector<std::vector<double>> tab(m + 2, std::vector<double>(n + 2));
  for (int i = 0; i < m; i++) {
    for (int j = 0; j < n; j++) {
      tab[i][j] = a[i][j];
    }
    basis[i] = n + i;
    tab[i][n] = -1;
    tab[i][n + 1] = b[i];
  }
  for (int j = 0; j < n; j++) {
    nonbasis[j] = j;
    tab[m][j] = maximize ? -c[j] : c[j];
  }
  nonbasis[n] = -1;
  tab[m + 1][n] = 1;
  auto pivot = [&](int row, int col) {
    double inv = 1.0 / tab[row][col];
    for (int i = 0; i < m + 2; i++) {
      if (i == row || fabs(tab[i][col]) <= eps) {
        continue;
      }
      double factor = tab[i][col] * inv;
      for (int j = 0; j < n + 2; j++) {
        if (j != col) {
          tab[i][j] -= factor * tab[row][j];
        }
      }
      tab[i][col] = -factor;
    }
    for (int j = 0; j < n + 2; j++) {
      if (j != col) {
        tab[row][j] *= inv;
      }
    }
    tab[row][col] = inv;
    std::swap(basis[row], nonbasis[col]);
  };
  auto simplex = [&](int phase) {
    int objective = (phase == 1 ? m + 1 : m);
    while (true) {
      int col = -1;
      for (int j = 0; j <= n; j++) {
        if (phase == 2 && nonbasis[j] == -1) {
          continue;
        }
        if (col == -1 || tab[objective][j] < tab[objective][col] - eps ||
            (fabs(tab[objective][j] - tab[objective][col]) <= eps && nonbasis[j] < nonbasis[col])) {
          col = j;
        }
      }
      if (tab[objective][col] >= -eps) {
        return true;
      }
      int row = -1;
      for (int i = 0; i < m; i++) {
        if (tab[i][col] <= eps) {
          continue;
        }
        double cur = tab[i][n + 1] / tab[i][col];
        double best = row == -1 ? 0 : tab[row][n + 1] / tab[row][col];
        if (row == -1 || cur < best - eps || (fabs(cur - best) <= eps && basis[i] < basis[row])) {
          row = i;
        }
      }
      if (row == -1) {
        return false;
      }
      pivot(row, col);
    }
  };
  int row = 0;
  for (int i = 1; i < m; i++) {
    if (tab[i][n + 1] < tab[row][n + 1]) {
      row = i;
    }
  }
  if (m > 0 && tab[row][n + 1] < -eps) {
    pivot(row, n);
    if (!simplex(1) || fabs(tab[m + 1][n + 1]) > eps) {
      return -1;
    }
    for (int i = 0; i < m; i++) {
      if (basis[i] != -1) {
        continue;
      }
      int col = -1;
      for (int j = 0; j <= n; j++) {
        if (fabs(tab[i][j]) > eps && (col == -1 || nonbasis[j] < nonbasis[col])) {
          col = j;
        }
      }
      if (col != -1) {
        pivot(i, col);
      }
    }
  }
  if (!simplex(2)) {
    return 1;
  }
  x->assign(n, 0);
  for (int i = 0; i < m; i++) {
    if (basis[i] < n) {
      (*x)[basis[i]] = tab[i][n + 1];
    }
  }
  return 0;
}
\end{lstlisting}
\Needspace*{4\baselineskip}
\textbf{Example Usage}\par
\begin{lstlisting}[language=C++]
#include <iostream>
using namespace std;

int main() {
  // Solve [x, y] that maximizes 3x + 4y, subject to x, y >= 0 and:
  //  -2x +    1y <=  0
  //   1x + 0.85y <=  9
  //   1x +    2y <= 14
  vector<vector<double>> va{{-2, 1}, {1, 0.85}, {1, 2}};
  vector<double> vb{0, 9, 14};
  vector<double> vc{3, 4};
  vector<double> x;
  assert(simplex_solve(va, vb, vc, &x) == 0);
  double maxval = 0;
  for (int i = 0; i < static_cast<int>(x.size()); i++) {
    maxval += vc[i] * x[i];
  }
  cout << "Solution = " << maxval << " at (" << x[0];
  for (int i = 1; i < static_cast<int>(x.size()); i++) {
    cout << ", " << x[i];
  }
  cout << ")." << endl;

  // Feasible even though the first RHS is negative: x >= 1 and x <= 3.
  vector<vector<double>> negative_rhs{{-1}, {1}};
  vector<double> negative_bounds{-1, 3};
  vector<double> one_var{1};
  assert(simplex_solve(negative_rhs, negative_bounds, one_var, &x) == 0);
  assert(fabs(x[0] - 3) < 1e-9);

  // Infeasible: x <= -1 contradicts x >= 0.
  vector<vector<double>> infeasible{{1}};
  vector<double> bad_bounds{-1};
  assert(simplex_solve(infeasible, bad_bounds, one_var, &x) == -1);

  // Unbounded: maximize x subject only to x >= 0.
  vector<vector<double>> unbounded{{-1}};
  vector<double> lower_bound{0};
  assert(simplex_solve(unbounded, lower_bound, one_var, &x) == 1);
  return 0;
}
\end{lstlisting}
\begin{exampleoutput}
Solution = 33.3043 at (5.30435, 4.34783).
\end{exampleoutput}
